\section{Antecedentes}

Fortalecer las habilidades blandas en los ingenieros de sistemas en formación es fundamental para su desempeño y proyección en el entorno laboral. En coherencia con esta necesidad, diversas investigaciones y trabajos previos han identificado las competencias más relevantes, analizado su efecto en el desempeño laboral y propuesto lineamientos para su integración en la formación académica.

\subsection{A Systematic Mapping Study on Soft Skills in Software Engineering}

Desde un enfoque investigativo, la revista académica de Universal Computer Science, en su publicación \textbf{“A Systematic Mapping Study on Soft Skills in Software Engineering”}, escrita por Gerardo Matturo, Florencia Raschetti y Carina Fontán, informa sobre los resultados de un estudio de mapeo sistemático destinado a identificar el estado del arte sobre habilidades blandas en ingeniería de software y a determinar cuáles habilidades son relevantes para su práctica, basándose en la frecuencia con la que son mencionadas en la literatura revisada por los autores. Su objetivo es ayudar a los gerentes de Recursos Humanos y líderes de equipos de proyectos de software a conocer qué habilidades no técnicas son generalmente demandadas por la industria e identificar cuáles de ellas pueden ser adecuadas para su entorno \cite{Matturro2019}.

\subsection{Skills Development for Software Engineers: Systematic Literature Review}

De manera complementaria, en el artículo \textbf{“Skills development for software engineers: Systematic literature review”} se identifican las habilidades blandas necesarias para los futuros ingenieros de software y las metodologías de enseñanza que contribuyen a desarrollar tales habilidades desde el pregrado, manteniendo a los estudiantes motivados. Además, se propone un marco de trabajo para ayudar a los educadores a llevar a cabo un proceso de enseñanza-aprendizaje que incluya habilidades técnicas y blandas durante el pregrado de futuros ingenieros de software \cite{Borges2024}.

\subsection{Acquiring Professional Software Engineering Skills through Studio-based Learning}

La 17ª Conferencia Internacional sobre Educación Superior y Capacitación Basada en Tecnología de la Información (ITHET) 2018 presentó un estudio detallado sobre la implementación del aprendizaje basado en estudio (SBL) en el programa de Ingeniería de Software de la Universidad de Monmouth, titulado \textbf{“Acquiring Professional Software Engineering Skills through Studio-based Learning”} \cite{Rosca2018}. Esta implementación se centra en preparar a los estudiantes para ser contribuyentes exitosos en sus futuras empresas, inscribiéndolos en una experiencia SBL de dos semestres durante su último año. En este curso, los estudiantes trabajan en proyectos de software que simulan el ciclo de vida del desarrollo de software en un entorno real y practican actividades específicas de la industria que desarrollan habilidades profesionales clave, como el trabajo en equipo, la comunicación efectiva, el aprendizaje continuo y la capacidad de funcionar en equipos interdisciplinarios. El estudio no solo detalla cómo se ha llevado a cabo esta implementación, sino que también la compara con el marco definido por los investigadores de la Universidad de Lancaster. Una característica innovadora de este enfoque es que introduce un proceso de aprendizaje casi completamente dirigido por los estudiantes, con solo un marco de referencia para guiarlos, lo cual ha demostrado ser eficaz en el desarrollo de habilidades profesionales en los estudiantes.

\subsection{Soft Skills: The Software Developer's Life Manual}

En el ámbito editorial, el libro escrito por John Sonmez, titulado \textbf{“Soft Skills: The Software Developer’s Life Manual”}, es una guía práctica centrada en las habilidades blandas y en el bienestar general del desarrollador de software. Sonmez aborda una amplia gama de temas importantes, organizados en siete secciones: carrera y productividad, comercialización personal, aprendizaje y educación continua, habilidades de comunicación, finanzas personales, salud y estilo de vida y, por último, desarrollo personal. En cada sección brinda ejemplos de su propia experiencia y propone actividades para que los lectores puedan poner en práctica lo aprendido \cite{Sonmez2014}.

\subsection{Learning How to Learn: Aprendiendo a Aprender}

El curso \textbf{“Aprendiendo a aprender: Poderosas herramientas mentales con las que podrás dominar temas difíciles (Learning How to Learn)”} \cite{coursera2024} ofrece acceso a técnicas de aprendizaje invaluables utilizadas por expertos en diversas disciplinas, como arte, música, literatura, matemáticas, ciencia, deportes, y muchas más. Con estos enfoques, los estudiantes, independientemente de su nivel de habilidad en los temas que deseen dominar, pueden transformar su forma de pensar y su vida. Aquellos con experiencia avanzada obtienen una visión profunda de la maquinaria cerebral, con ideas para fomentar un nuevo y poderoso aprendizaje, consejos contraintuitivos para presentar exámenes y perspectivas para optimizar el tiempo dedicado a tareas y resolución de problemas. Aquellos que enfrentan dificultades descubren un conjunto de técnicas prácticas estructuradas para mejorar su desempeño.

\subsection{Desarrollo de una Herramienta de Apoyo para el Aprendizaje de Habilidades Blandas con Gamificación}

En el ámbito institucional, varios trabajos de grado desarrollados en la Universidad del Valle se articulan con la necesidad de fortalecer las habilidades blandas mediante propuestas tecnológicas y enfoques basados en gamificación. Entre ellos se encuentra el trabajo de Daniel Felipe Cossio Marulanda, titulado \textbf{“Desarrollo de una herramienta de apoyo para el aprendizaje de habilidades blandas con gamificación”} \cite{cossioMarulanda2024}, cuyo objetivo fue construir un prototipo o aplicación semilla orientada al desarrollo de habilidades blandas para el entorno laboral, con énfasis en el trabajo en equipo y el liderazgo. Este proyecto sirve como base para la implementación de diversos módulos que faciliten la práctica y adquisición de habilidades interpersonales a menudo pasadas por alto en la formación académica tradicional.

\subsection{Desarrollo de Módulos de Mindset y Fitness para Plataforma Gamificada}

Juan Camilo Azcárate Cárdenas, en su trabajo de grado \textbf{“Desarrollo de módulos de mindset y fitness para el prototipo de plataforma enfocada en la enseñanza de habilidades blandas mediante gamificación”} \cite{azcarateCardenas2024}, propone integrar dos nuevos módulos específicos: “Fitness”, orientado al bienestar general, y “Mindset”, destinado a transformar las reacciones ante la vida de reactivas a proactivas. Estos módulos buscan complementar y enriquecer la experiencia de aprendizaje de habilidades blandas para estudiantes y profesionales de ingeniería de sistemas.

\subsection{Creación de los Módulos de Productividad y Finanzas para Plataforma Gamificada}

El trabajo de grado de Daniel Gómez Suárez, titulado \textbf{“Creación de los módulos de productividad y finanzas para prototipo de plataforma web enfocada al aprendizaje de habilidades blandas”} \cite{gomezSuarez2024}, se centra en determinar la estrategia más adecuada para desarrollar e integrar los módulos de productividad y finanzas en la plataforma, orientada al fortalecimiento de habilidades blandas en estudiantes y profesionales de ingeniería de sistemas.

\subsection{Creación de Módulos Interactivos para el Desarrollo Continuo de la Oratoria y el Pensamiento Crítico}

El trabajo de grado titulado \textbf{“Creación de módulos interactivos para el Desarrollo Continuo de la Oratoria y el Pensamiento Crítico en el prototipo de enseñanza de habilidades blandas”}, desarrollado por Juan David Moreno Alfonso \cite{morenoAlfonso2025}, propone el diseño e implementación de dos módulos adicionales para la plataforma institucional orientada al fortalecimiento de habilidades blandas.
\\
\\
El proyecto se centra en las competencias de \textit{oratoria} y \textit{pensamiento crítico}, consideradas fundamentales para el desempeño académico y profesional. Para ello, el autor desarrolla módulos interactivos que permiten a los estudiantes practicar habilidades comunicativas y argumentativas mediante actividades guiadas, retroalimentación continua y ejercicios diseñados para simular escenarios reales de debate y exposición oral. 
\\
\\
El trabajo también documenta la arquitectura de software utilizada, el flujo de interacción entre pantallas, la infraestructura del sistema y la integración del backend y frontend necesarios para el funcionamiento de los módulos. Este proyecto constituye un antecedente relevante, pues amplía el alcance del prototipo institucional de habilidades blandas mediante nuevos componentes que fortalecen las capacidades comunicativas y cognitivas de los estudiantes, alineándose con los objetivos formativos de la Escuela de Ingeniería de Sistemas y Computación de la Universidad del Valle.

